\chapter{Aritmetika Floating-Point dan Analisis Galat}

Sebuah program menyelesaikan $x^{2}+10^{8}x+1=0$ dengan rumus kuadrat yang kita pakai
sejak sekolah, $x=(-b\pm\sqrt{b^{2}-4ac})/(2a)$. Untuk $a=1$, $b=10^{8}$, $c=1$ ia
melaporkan akar $-10^{8}$ dan $-7{,}450580596923828\times10^{-9}$. Akar pertama benar
sampai digit terakhir. Akar kedua salah seperempat: nilai sebenarnya $-10^{-8}$, dan
galat relatifnya $0{,}255$. Rumusnya tidak salah tulis. Mesinnya juga tidak rusak.
Koefisiennya tepat, diskriminannya tepat, dan pengurangannya bahkan tepat menurut
teorema yang akan kita buktikan di bawah.

Yang runtuh adalah bentuk aljabarnya. Bilangan $b=10^{8}$ dan $\sqrt{b^{2}-4ac}$ berbeda
hanya satu ulp, yaitu satu langkah terkecil yang masih bisa dibedakan mesin pada skala
itu: selisihnya $1{,}49\times10^{-8}$, atau $1{,}5\times10^{-16}$ kali besar bilangannya
sendiri. Selisih sekecil itu tidak berarti apa-apa selama kita memandang bilangan
berukuran $10^{8}$. Begitu keduanya dikurangkan, hasilnya berukuran $10^{-8}$ dan
selisih tadi menjadi seluruh jawabannya: yang tersisa hanyalah galat, bukan angka.

Bab 4 dan Bab 5 membangun limit, kekontinuan, dan turunan di atas $\mathbb{R}$ yang
setiap titiknya eksak. Bab 20 memakai rumus variansi dan akar kuadrat dengan anggapan
yang sama. Komputer tidak bekerja di $\mathbb{R}$. Ia bekerja di himpunan berhingga
bilangan yang punya aturan sendiri: setiap operasi dibulatkan, penjumlahan tidak
asosiatif, dan dua cara aljabar yang setara bisa memberi jawaban berbeda jauh. Bab ini
membangun cara membaca aturan itu, mengukur kerusakannya sebelum program dijalankan,
dan memilih bentuk rumus yang tidak runtuh.

\begin{untukanak}
Bayangkan sebuah timbangan dapur digital yang layarnya hanya muat lima angka. Kita
menimbang satu karung beras bersama wadahnya, lalu menimbang wadah kosongnya, lalu
mengurangkan kedua angka itu untuk mendapat berat berasnya. Kalau wadah dan beras
sama-sama beratnya sekitar dua puluh kilogram, layar hanya bisa menunjukkan angka
sampai puluhan gram terdekat. Selisih dua angka besar yang hampir sama itu memang
benar secara hitungan, tetapi angka terakhirnya sudah tidak berarti apa-apa. Beras
yang sebenarnya delapan ratus gram bisa muncul sebagai tujuh ratus gram, atau bahkan
angka yang tidak masuk akal.

Ada masalah kedua yang lebih halus. Timbangan yang sama tidak bisa menimbang satu
butir beras: jarumnya tidak bergerak sama sekali. Kalau kita menuang sejuta butir satu
per satu, tidak ada satu pun butir yang tercatat, padahal sejuta butir itu berat.
Mesin hitung bekerja persis seperti itu: ada bilangan yang terlalu kecil untuk
mengubah apa pun, dan ada pengurangan yang membuang hampir semua angka.

Pertanyaannya: kalau mesinnya memang begitu, bagaimana kita bisa tahu kapan jawabannya
masih bisa dipercaya? Di bab ini kita akan tahu mengapa kesalahan itu bukan kejadian
acak, melainkan sesuatu yang bisa dihitung batasnya sebelum satu baris program pun
dijalankan.
\end{untukanak}

\section{Bilangan yang tidak punya wakil}
\tingkat{dasar}

Di antara $0$ dan $1$ ada tak berhingga banyak bilangan real. Di dalam komputer, semua
bilangan disimpan dalam 32 atau 64 bit, jadi yang tersedia hanya berhingga banyak.
Kompromi yang diambil sejak mesin pertama sampai standar IEEE-754 adalah notasi ilmiah
biner dengan lebar bidang tetap: satu tanda, satu eksponen, satu mantissa.

\kamus{sistem floating-point}{floating-point system}{Notasi ilmiah biner dengan jumlah
digit tetap, seperti menulis $1{,}01\times2^{7}$ dengan tempat terbatas.}

\begin{definisi}[Sistem bilangan floating-point]
Misalkan $\beta=2$ basis, $p$ presisi, dan $e_{\min}\le e\le e_{\max}$ batas eksponen.
Bilangan normal adalah
\[
  x=\pm\left(1.d_{1}d_{2}\dots d_{p-1}\right)_{2}\times2^{e},
\]
dengan $d_{i}\in\{0,1\}$. Bilangan subnormal memakai $e=e_{\min}$ dengan digit
terdepan $0$.\index{subnormal}\index{aritmetika floating-point} Himpunan seluruh
bilangan yang terwakili, ditambah nol, kedua tak hingga, dan NaN, kita sebut
$\mathbb{F}$.
\end{definisi}

Digit terdepan selalu $1$ untuk bilangan normal, jadi digit itu tidak disimpan;
itulah sebabnya bidang mantissa punya $p-1$ bit. Satu bit yang dihemat menggandakan
presisi: 23 bit tersimpan memberi 24 bit
signifikan.



\begin{teorema}[Banyaknya bilangan mesin]
Untuk setiap eksponen $e$, banyaknya bilangan normal dalam selang
$\left[2^{e},2^{e+1}\right)$ adalah $2^{p-1}$. Himpunan $\mathbb{F}$ berhingga
banyaknya, jadi terhitung.
\end{teorema}

\begin{proof}
Pada selang itu setiap bilangan berbentuk $M\cdot2^{e-p+1}$ dengan $M$ bilangan bulat
dan $2^{p-1}\le M<2^{p}$. Banyaknya $M$ semacam itu tepat $2^{p}$, dikurangi
$2^{p-1}$, yaitu $2^{p-1}$. Ada berhingga banyak eksponen yang diizinkan, jadi
$\mathbb{F}$ adalah gabungan berhingga himpunan berhingga.
\end{proof}

Untuk binary64, $2^{p-1}=2^{52}=4\,503\,599\,627\,370\,496$. Itu banyak, tetapi Bab 2
sudah menunjukkan $\mathbb{R}$ tak terhitung. Jadi hampir setiap bilangan real tidak
punya wakil: representasi selalu berarti pembulatan, dan pembulatan selalu berarti
galat. Pertanyaannya bukan apakah ada galat, melainkan seberapa besar dan ke mana ia
pergi.

\begin{bacakeras}
Dua suara, dibaca bergantian dengan teman. Ketuk meja setiap kali kata ``potongan''
diucapkan.

``Satu dibagi tiga, berapa?'' ``Nol koma tiga tiga tiga, terus, tanpa henti.''
``Mesin menulisnya sampai habis?'' ``Tidak. Ia berhenti di tempat yang muat.''
``Jadi yang disimpan apa?'' ``Potongan. Bukan nilainya, potongannya.''
``Kalau potongannya dijumlahkan?'' ``Yang bertambah potongan, bukan nilai.''
``Kalau potongannya dikurangkan?'' ``Angka depannya habis, yang tinggal debu.''

Ulangi tiga baris tengah, makin cepat, dengan ketukan yang sama: potongan bukan nilai,
potongan bukan nilai, potongan bukan nilai.

Suara kedua menutup pelan: mesin tidak berdusta, ia hanya memotong. Yang keliru adalah
kita kalau membaca potongan sebagai nilai.

Baris yang harus diulang pembaca, sekali, dengan suara jelas: komputer menyimpan
potongan, bukan nilai, jadi setiap hasil hitung harus diperiksa, bukan dipercaya.
\end{bacakeras}

Standar IEEE-754 menetapkan dua lebar yang dipakai di mana-mana. Tabel ini memuat
seluruh angkanya; tidak ada yang perlu dihafal kecuali $p$ dan biasnya.

\begin{center}
\begin{tabular}{lcc}
\hline
 & binary32 & binary64 \\
\hline
bit tanda / eksponen / mantissa & $1/8/23$ & $1/11/52$ \\
presisi $p$ & $24$ & $53$ \\
bias eksponen & $127$ & $1023$ \\
$e_{\min},\ e_{\max}$ & $-126,\ 127$ & $-1022,\ 1023$ \\
$\mathrm{eps}=2^{-p+1}$ & $1{,}1920929\times10^{-7}$ & $2{,}2204460\times10^{-16}$ \\
$u=2^{-p}$ & $5{,}9604645\times10^{-8}$ & $1{,}1102230\times10^{-16}$ \\
terkecil normal & $1{,}1754944\times10^{-38}$ & $2{,}2250739\times10^{-308}$ \\
terkecil subnormal & $1{,}4012985\times10^{-45}$ & $4{,}9406565\times10^{-324}$ \\
terbesar & $3{,}4028235\times10^{38}$ & $1{,}7976931\times10^{308}$ \\
\hline
\end{tabular}
\end{center}

Eksponen disimpan dengan bias supaya tidak perlu tanda sendiri: nilai tersimpan $123$
berarti eksponen sebenarnya $123-127=-4$. Trik ini membuat urutan bit dari nol sampai
tak hingga positif naik monoton, sehingga dua bilangan positif dapat dibandingkan
sebagai bilangan bulat.

Mari bongkar tiga pola bit yang benar-benar dihasilkan Python. Pada binary32,
$0{,}1$ tercetak sebagai \texttt{0 01111011 10011001100110011001101} dengan
$e=123-127=-4$; $6$ menjadi \texttt{0 10000001 10000000000000000000000}, yaitu
$1{,}5\times2^{2}$; dan $-12{,}375$ menjadi
\texttt{1 10000010 10001100000000000000000}, yaitu $-1{,}546875\times2^{3}$.
Dua kasus batas melengkapi gambarannya:

\begin{center}
\begin{tabular}{ll}
\hline
subnormal terkecil & \texttt{0 00000000 00000000000000000000001} \\
terbesar & \texttt{0 11111110 11111111111111111111111} \\
\hline
\end{tabular}
\end{center}

Yang pertama bernilai $2^{-149}$: seluruh bit eksponennya nol dan hanya bit terakhir
mantissa menyala. Yang kedua bernilai $3{,}4028235\times10^{38}$, dengan eksponen
tersimpan $254$ dan mantissa penuh.

Bentuk yang lebih mudah dibaca untuk binary64 disediakan oleh \texttt{float.hex()}. Hasilnya untuk beberapa bilangan:

\begin{center}
\begin{tabular}{ll}
\hline
$0{,}1$ & \texttt{0x1.999999999999ap-4} \\
$6{,}0$ & \texttt{0x1.8000000000000p+2} \\
$10^{-8}$ & \texttt{0x1.5798ee2308c3ap-27} \\
$5\times10^{-324}$ & \texttt{0x0.0000000000001p-1022} \\
\hline
\end{tabular}
\end{center}

Bentuk itu memperlihatkan tiga hal. Digit di belakang titik heksadesimal selalu tiga
belas, dan $13\times4=52$, yaitu lebar mantissa binary64. Digit terakhir $0{,}1$
adalah \texttt{a}, yaitu $1010$, hasil pembulatan ke atas dari pola $1001$ yang
berulang. Baris terakhir punya nol di depan titik: tanda bilangan subnormal.

Nilai khusus melengkapi sistem ini: $\pm\infty$ untuk hasil yang melampaui
$e_{\max}$, nol bertanda, bilangan subnormal di bawah $2^{e_{\min}}$, dan NaN
(\emph{not a number}) untuk bentuk tak tentu seperti $0/0$. Satu sifat NaN layak
diingat karena sering mengejutkan: ia tidak sama dengan dirinya sendiri, sehingga
perbandingan \texttt{x == x} bernilai salah justru ketika $x$ adalah NaN.\index{NaN}

\section{Satu-satunya jembatan: pembulatan}
\tingkat{dasar}

Tidak ada operasi lain yang menghubungkan $\mathbb{R}$ dengan $\mathbb{F}$ selain
pembulatan. Karena itu seluruh analisis galat berdiri di atas satu pertanyaan: seberapa
jauh bilangan mesin terdekat dari bilangan aslinya, dan bagaimana jarak itu berskala.

IEEE-754 menetapkan aturan \emph{round to nearest, ties to even}: pilih bilangan mesin
terdekat, dan bila dua kandidat sama jauhnya pilih yang digit terakhir mantissanya
genap. Aturan kedua bukan kemewahan teoretis. Tanpa aturan itu, pembulatan pada titik
seri selalu condong ke satu arah, dan galat sejuta operasi akan menumpuk ke arah yang
sama, bukan saling meniadakan.

Titik seri itu benar-benar muncul, dan kita bisa membuatnya dengan sengaja. Dalam
binary64, jarak antara $1$ dan bilangan sesudahnya adalah $2^{-52}$. Maka
$1+2^{-53}$ jatuh tepat di tengah, dan aturan genap memilih $1$ karena mantissanya
berakhiran nol. Python memberi

\begin{center}
\texttt{1 + 2.0**-53 = 1.0}, sedangkan \texttt{1 + 3*2.0**-53 = 1.0000000000000004}.
\end{center}

Kasus pertama membulat ke bawah, kasus kedua ke atas, padahal keduanya sama-sama tepat
setengah jalan. Yang membedakan hanya paritas digit terakhir. Inilah sebabnya
$1+10^{-16}=1$ dalam binary64: bukan karena $10^{-16}$ diabaikan dengan sembarangan,
tetapi karena tidak ada tempat untuk menyimpannya.

\kamus{epsilon mesin}{machine epsilon}{Jarak antara $1$ dan bilangan mesin
berikutnya: ukuran satu langkah di skala satu.}\index{epsilon mesin}

Dua angka sering tertukar dan pertukarannya mengubah batas galat dengan faktor dua.
Epsilon mesin $\mathrm{eps}=2^{-p+1}$ adalah jarak antara $1$ dan bilangan
sesudahnya, yaitu $2^{-52}$ untuk binary64. Satuan pembulatan
$u=2^{-p}=2^{-53}$ adalah setengah jarak itu, dan hanya $u$ yang muncul dalam batas
galat. Python menyebut yang pertama \texttt{numpy.finfo(float).eps}; yang kedua harus
kita hitung sendiri.

Jarak antar bilangan tidak seragam. Di dekat $1$ jaraknya $2^{-52}$, di dekat
$10^{10}$ jaraknya $1{,}9\times10^{-6}$. Yang seragam adalah jarak itu relatif
terhadap besar bilangannya, dan itulah isi lemma berikut.

\begin{lemma}[Ulp]
Jika $2^{e}\le|x|<2^{e+1}$, maka jarak antara $x$ dan bilangan mesin berikutnya adalah
$\operatorname{ulp}(x)=2^{e-p+1}$.
\end{lemma}

\begin{proof}
Pada selang itu seluruh bilangan normal berbentuk $M\cdot2^{e-p+1}$ dengan $M$ bulat
dan $2^{p-1}\le M<2^{p}$; banyaknya $M$ adalah $2^{p-1}$ dan selisih dua $M$
berurutan adalah satu satuan $2^{e-p+1}$. Selang berikutnya memakai $e+1$ sehingga
jaraknya berlipat dua, dan itulah satu-satunya titik tempat kerapatan berubah.
\end{proof}

\kamus{ulp}{unit in the last place}{Satu langkah terkecil yang bisa dibedakan mesin
pada skala sebuah bilangan.}\index{ulp}

Lemma itu menjelaskan dua hal sekaligus.
$\operatorname{ulp}(1)=2^{-52}=2{,}220446049250313\times10^{-16}$, sedangkan
$\operatorname{ulp}(10^{10})=2^{-52}\cdot2^{33}\approx1{,}9\times10^{-6}$. Jadi
menambahkan $1$ kepada $10^{16}$ tidak mengubah apa pun dalam binary64, karena
$\operatorname{ulp}(10^{16})=2$: memang tidak ada bilangan mesin di antara keduanya.

Gambar~\ref{fig:titikmesin} memperlihatkan akibat lemma itu pada empat oktaf
pertama: setiap oktaf memuat $2^{23}$ titik, tetapi lebarnya berlipat dua, sehingga
jarak antar titik berlipat dua pula.

\begin{figure}[htbp]
\centering
\begin{tikzpicture}[x=2.1cm,y=1cm,every node/.style={font=\footnotesize}]
  \draw[ink,thick,->] (-0.55,0) -- (4.30,0);
  \draw[ink,thick,fill=white] (-0.42,0) circle (1.7pt);
  \draw[ink,thick] (-0.48,-0.09) -- (-0.36,0.09);
  \foreach \x in {0,0.125,...,0.875} {\draw[ink] (\x,0.055) -- (\x,-0.055);}
  \foreach \x in {1,1.25,...,1.875} {\draw[ink] (\x,0.055) -- (\x,-0.055);}
  \foreach \x in {2,2.5,...,3.0} {\draw[ink] (\x,0.055) -- (\x,-0.055);}
  \foreach \x in {3,4} {\draw[ink] (\x,0.055) -- (\x,-0.055);}
  \foreach \x/\lab in {0/{$2^{0}$},1/{$2^{1}$},2/{$2^{2}$},3/{$2^{3}$},4/{$2^{4}$}} {
    \draw[ink,thick] (\x,0.14) -- (\x,-0.14);
    \node[below=1pt] at (\x,-0.14) {\lab};
  }
  \draw[{Latex}-{Latex},accent] (0.125,-0.62) -- (0.25,-0.62) node[midway,below] {$d$};
  \draw[{Latex}-{Latex},accent] (1.25,-0.62) -- (1.50,-0.62) node[midway,below] {$2d$};
  \draw[{Latex}-{Latex},accent] (2.50,-0.62) -- (3.00,-0.62) node[midway,below] {$4d$};
  \draw[{Latex}-{Latex},accent] (3.00,-0.62) -- (4.00,-0.62) node[midway,below] {$8d$};
  \node[align=center,text width=3.0cm] at (1.95,0.92)
    {tiap oktaf berisi $2^{23}$ titik, lebarnya dua kali oktaf sebelumnya};
  \node[align=center,text width=3.0cm] at (0.40,-1.32)
    {ke kiri menuju nol: jarak mengecil};
\end{tikzpicture}
\caption{Empat oktaf pertama bilangan mesin binary32. Setiap oktaf memuat $2^{23}$
titik, tetapi lebarnya dua kali oktaf sebelumnya, jadi jarak antar titik berlipat dua:
$d=2^{-23}$ pada $[1,2)$, $2d$ pada $[2,4)$, $4d$ pada $[4,8)$, dan $8d$ pada
$[8,16)$. Titik di ekor tidak lebih sedikit, hanya lebih jarang; ke arah nol jaraknya
mengecil sampai batas subnormal.}
\label{fig:titikmesin}
\end{figure}

Ada satu kehalusan yang perlu disebut sebelum kita memakai ulp sebagai satuan galat.
Galat pembulatan paling besar setengah ulp, tetapi bila dinyatakan sebagai galat
relatif, angka itu bukan satu bilangan melainkan satu selang. Untuk $x$ yang tepat
sama dengan pangkat dua, setengah ulp relatifnya $2^{-53}$; untuk $x$ yang hampir
mencapai pangkat dua berikutnya, nilainya mendekati $2^{-54}$. Jadi galat relatif satu
pembulatan bisa berayun dengan faktor dua. Goldberg menamai ayunan itu \emph{wobble}, dan
karena itu batas $u$ selalu dipakai sebagai batas terburuk, bukan sebagai nilai khas.

\begin{teorema}[Galat pembulatan relatif]
Untuk setiap $x\in\mathbb{R}$ dalam jangkauan normal, pembulatan benar memenuhi
$\operatorname{fl}(x)=x(1+\delta)$ dengan $|\delta|\le u=2^{-p}$.
\end{teorema}

\begin{proof}
Ambil $2^{e}\le|x|<2^{e+1}$. Bilangan mesin terdekat berjarak paling jauh setengah
ulp, yaitu $2^{e-p}$, sehingga
$|\operatorname{fl}(x)-x|\le2^{e-p}=u\cdot2^{e}\le u|x|$. Membagi dengan $|x|$
memberi $|\delta|\le u$, dan tanda $\delta$ mengikuti arah pembulatan.
\end{proof}

Bentuk $x(1+\delta)$ adalah awal dari segalanya: ia mengubah pertanyaan ``seberapa
salah'' menjadi ``berapa banyak faktor $(1+\delta)$ yang menumpuk''. Operasi
$+,-,\times,\div,\sqrt{\phantom{x}}$ dalam IEEE-754 wajib mengembalikan pembulatan
benar dari hasil eksaknya, jadi setiap operasi menyumbang satu faktor
semacam itu dan tidak lebih. Jaminan inilah yang membuat analisis galat mungkin.


\section{Galat, bilangan kondisi, dan pembatalan}
\tingkat{menengah}

Sekarang kita punya satuan. Yang belum ada adalah ukuran seberapa besar galat sebuah
perhitungan dibandingkan galat masukannya. Dua masalah bisa sama-sama menerima
masukan yang salah $10^{-16}$ dan keluar dengan jawaban yang salah $10^{-16}$ atau
$10^{-1}$, bergantung pada masalahnya sendiri.

\kamus{bilangan kondisi}{condition number}{Pengali yang menunjukkan berapa kali galat
masukan diperbesar oleh sebuah perhitungan.}\index{bilangan kondisi}

\begin{definisi}[Galat absolut, relatif, dan ke belakang]
Misalkan $\tilde{x}$ penghampiran bagi $x$. Galat absolutnya $|\tilde{x}-x|$, galat
relatifnya $|\tilde{x}-x|/|x|$ untuk $x\neq0$. Jika $\tilde{x}$ adalah jawaban eksak
bagi masalah dengan data $x+\Delta x$, maka $\|\Delta x\|/\|x\|$ disebut galat ke
belakang, dan algoritma yang selalu memberi galat ke belakang kecil disebut
\emph{stabil ke belakang}.\index{galat relatif}\index{stabil ke belakang}
\end{definisi}

Pembedaan terakhir paling berguna dalam praktik. Algoritma stabil ke belakang boleh
salah jauh, asalkan jawabannya benar untuk masalah yang sedikit berbeda:
kesalahannya milik masalah. Galat ke belakang yang besar adalah kesalahan penulis
kode.

Untuk fungsi satu peubah, pembesaran itu dihitung langsung dengan turunan: bilangan
kondisi adalah kemiringan fungsi di titik itu, dikali besar titiknya, dibagi besar
nilainya.

\begin{teorema}[Bilangan kondisi]
Misalkan $f$ terturunkan pada selang yang memuat $x$ dan $f(x)\neq0$. Bila
$\tilde{x}=x(1+\delta)$ dengan $|\delta|\le u$, maka
\[
  \frac{|f(\tilde{x})-f(x)|}{|f(x)|}\;\le\;\kappa_{f}(x)\,u+O(u^{2}),
  \qquad \kappa_{f}(x)=\frac{|x\,f'(x)|}{|f(x)|}.
\]
\end{teorema}

\begin{proof}
Menurut teorema nilai rata-rata, ada $\xi$ di antara $x$ dan $\tilde{x}$ dengan
$f(\tilde{x})-f(x)=f'(\xi)(\tilde{x}-x)$. Karena $|\tilde{x}-x|\le u|x|$ dan $f'$
kontinu di $x$, berlaku $|f'(\xi)|=|f'(x)|+O(u)$. Maka
$|f(\tilde{x})-f(x)|\le u|x|\,|f'(x)|+O(u^{2}|x|)$. Membagi dengan $|f(x)|$ memberi
ketaksamaan yang diklaim.
\end{proof}

Bilangan $\kappa_{f}$ tidak bergantung pada mesin, hanya pada masalahnya. Untuk
$f(x)=\sqrt{x}$ ia bernilai $1/2$: akar kuadrat tidak pernah memperburuk galat relatif.
Untuk $f(x)=\log x$ ia bernilai $1/|\log x|$, yang besar ketika $x$ dekat $1$ karena
di situ $\log x$ hampir nol. Untuk $f(x)=x-c$ dengan $c$ konstanta besar, ia bernilai
$|x|/|x-c|$, yang membesar tanpa batas ketika $x$ mendekati $c$. Pola terakhir itulah
yang menjadi sumber seluruh persoalan bab ini.

Pengurangan dua bilangan yang berdekatan adalah operasi paling berbahaya dalam
aritmetika mesin, dan bahayanya sering salah alamat. Kesalahpahaman lazim berbunyi:
pengurangan kehilangan digit, jadi operasi pengurangan itu sendiri tidak akurat. Yang
benar lebih halus dan lebih menenteramkan: bila kedua bilangan berdekatan,
pengurangannya justru eksak.

\begin{teorema}[Sterbenz]
Jika $a,b\in\mathbb{F}$ dan $b/2\le a\le2b$ dengan $b>0$, maka $a-b$ terwakili persis
di $\mathbb{F}$.
\end{teorema}

\begin{proof}
Misalkan $2^{e}\le b<2^{e+1}$. Setiap bilangan mesin pada selang itu kelipatan bulat
$2^{e-p+1}$, jadi $a$ juga, sehingga $a-b$ kelipatan $2^{e-p+1}$. Dari
$b/2\le a\le2b$ kita peroleh $0\le a-b\le b/2<2^{e}$. Kelipatan $2^{e-p+1}$ yang
terletak di $[0,2^{e})$ berbentuk $M\cdot2^{e-p+1}$ dengan $0\le M<2^{p}$, dan
semuanya terwakili dengan $p$ digit. Jadi hasilnya eksak.
\end{proof}

Jadi bila hasil pengurangan tampak buruk, yang buruk adalah masukannya. Angka yang
``hilang'' sudah hilang sebelum pengurangan dilakukan: $a$ dan $b$ masing-masing hanya
diketahui sampai galat relatif $u$, dan galat mutlak sebesar $u|a|$ itu tidak ikut
menyusut ketika kedua bilangan saling meniadakan. Yang menyusut hanya besar hasilnya,
sehingga galat relatifnya meledak.

\begin{proposisi}[Bilangan kondisi pengurangan]
Misalkan $x=a-b\neq0$ dan mesin bekerja dengan $\tilde{a}=a(1+\delta_{a})$,
$\tilde{b}=b(1+\delta_{b})$, $|\delta_{a}|,|\delta_{b}|\le u$. Maka
\[
  \frac{|\operatorname{fl}(\tilde{a}-\tilde{b})-x|}{|x|}
  \;\le\;\kappa\,u+u,
  \qquad \kappa=\frac{|a|+|b|}{|a-b|}.
\]
\end{proposisi}

\begin{proof}
Karena $\operatorname{fl}(\tilde{a}-\tilde{b})$ adalah pembulatan benar dari
$\tilde{a}-\tilde{b}$, galat tambahannya paling banyak $u|\tilde{a}-\tilde{b}|$.
Selanjutnya
$|\tilde{a}-\tilde{b}-(a-b)|\le|a||\delta_{a}|+|b||\delta_{b}|\le u(|a|+|b|)$,
dan membagi dengan $|x|=|a-b|$ memberi $u\kappa$. Menjumlahkan kedua galat dengan
ketaksamaan segitiga menyelesaikan bukti.
\end{proof}

Bilangan $\kappa$ tidak pernah lebih kecil dari satu dan membesar tanpa batas ketika
$|a-b|$ mengecil. Untuk $a=12345{,}678$ dan $b=12345{,}671$ kita peroleh $a-b=0{,}007$
dan $\kappa\approx3{,}53\times10^{6}$, sehingga galat relatif hasil bisa mencapai
$3{,}9\times10^{-10}$ walaupun setiap masukan hanya menyimpang $10^{-16}$. Karena setiap
faktor sepuluh dalam $\kappa$ menghapus satu digit signifikan,
$\log_{10}\kappa\approx6{,}5$ berarti enam setengah dari enam belas digit mesin hilang,
dan hasil $0{,}007$ hanya menyisakan sekitar sepuluh digit yang bisa dipercaya.

Satu catatan yang menyelamatkan: aturan Sterbenz berlaku untuk selisih, bukan untuk
penjumlahan. Untuk penjumlahan $a+b$ bilangan kondisi adalah
$(|a|+|b|)/|a+b|$, yang mendekati satu ketika $a$ dan $b$ bertanda sama. Menjumlahkan
bilangan besar tidak berbahaya. Mengurangkan dua bilangan besar yang hampir sama
berbahaya, dan itulah satu-satunya kebiasaan yang perlu diwaspadai.

\section{Pembatalan katastrofik: dua wajah}
\tingkat{menengah}

Kita sudah tahu mekanismenya. Sekarang kita beri namanya, karena dua jenis pembatalan
yang mirip sering dianggap sama padahal akibatnya berbeda jauh.

\begin{untukanak}
Kalau kita ingin tahu tebal selembar kertas, penggaris biasa tidak akan menolong: satu
lembar jauh lebih tipis daripada satu milimeter, jadi jarumnya seolah tidak bergerak.
Caranya bukan mencari penggaris yang lebih teliti, melainkan mengubah apa yang diukur:
tumpuk lima ratus lembar, ukur tebal tumpukannya, lalu bagi lima ratus. Hasilnya jauh
lebih dapat dipercaya, walaupun penggarisnya tetap penggaris yang sama.

Ternyata rumus juga bisa diperlakukan seperti itu. Ada rumus yang bekerja dengan
mengurangkan dua angka besar yang hampir sama, dan itu sama saja dengan menimbang
satu lembar kertas dengan penggaris. Ada rumus lain yang secara matematika sama
persis, tetapi menyusun langkahnya supaya angka kecil diukur langsung, bukan dicari
dari selisih. Hasilnya bisa berbeda jauh, padahal kertasnya tidak berubah.

Pertanyaannya: kalau dua rumus yang sama-sama benar memberi jawaban berbeda, mana yang
harus kita percaya? Di bab ini kita akan tahu mengapa rumus yang salah bukan yang
hitungannya keliru, melainkan yang memilih langkah yang membuang keterangan, dan
bagaimana mengenalinya sebelum program dijalankan.
\end{untukanak}

\kamus{pembatalan katastrofik}{catastrophic cancellation}{Hilangnya digit yang berarti
ketika dua bilangan hampir sama dikurangkan, sehingga yang tersisa hanya galat.}
\index{pembatalan katastrofik}

\begin{definisi}[Pembatalan katastrofik dan jinak]
Misalkan $x$ dan $y$ adalah bilangan mesin yang dihitung dari data sebelumnya.
Pengurangan $x-y$ disebut \emph{pembatalan jinak} bila $x$ dan $y$ tepat, sehingga
hasilnya eksak menurut teorema Sterbenz. Ia disebut \emph{pembatalan katastrofik} bila
$x$ dan $y$ sendiri sudah membawa galat pembulatan, sehingga digit yang sama-sama
mereka miliki ikut hilang dan yang tersisa didominasi galat itu.
\end{definisi}

Contoh jinak: $1{,}0000000000000002-1$ dalam binary64. Keduanya bilangan mesin,
selisihnya $2^{-52}$ dan terwakili persis. Hasilnya kecil, tetapi tidak ada satu digit
pun yang hilang: kita memang tahu kedua masukannya sampai bit terakhir.

Contoh katastrofik adalah rumus kuadrat di awal bab. Di situ $b^{2}$ dan $4ac$
masing-masing hasil perkalian, jadi keduanya dibulatkan. Untuk $a=1$, $b=10^{8}$,
$c=1$ keduanya kebetulan tepat, dan justru kekhususan itulah yang membuat contoh ini
layak diperiksa sampai bit: keruntuhannya bukan berasal dari perkalian, melainkan dari
pembulatan akar kuadratnya. Python menghitung

\begin{center}
\begin{tabular}{ll}
\hline
diskriminan $b^{2}-4ac$ & \texttt{9999999999999996.0} \\
$\sqrt{b^{2}-4ac}$ & \texttt{99999999.99999999} \\
$b-\sqrt{b^{2}-4ac}$ & \texttt{1.4901161193847656e-08} \\
$\operatorname{ulp}(b)$ & \texttt{1.4901161193847656e-08} \\
\hline
\end{tabular}
\end{center}

Nilai sebenarnya $b-\sqrt{b^{2}-4ac}$ adalah $2\times10^{-8}$, yaitu $1{,}34$ ulp.
Yang terhitung tepat $1$ ulp, dan galat setengah ulp pada akar kuadrat sudah cukup
untuk itu: akar kuadratnya benar sampai $0{,}46$ ulp pada skala $10^{8}$, tetapi
selisih yang seharusnya $1{,}34$ ulp hanya bisa mendarat pada $1$ atau $2$ ulp, jadi
informasinya tinggal satu pilihan biner. Hasil akhirnya $-7{,}45\times10^{-9}$, yaitu
$-2^{-27}$, sedangkan akar sebenarnya $-10^{-8}$. Galat relatif $0{,}255$.

Perhatikan pembagian tugasnya. Akar kuadrat bekerja dengan galat relatif
$5\times10^{-17}$, jauh di bawah $u$. Masalahnya sendiri berkondisi baik: menggeser
$b$ sebesar satu ulp hanya menggeser akar kecil sebesar satu ulp, jadi
$\kappa\approx|b|/\sqrt{b^{2}-4ac}\approx1$. Yang buruk adalah \emph{bentuk} yang kita
pilih: bilangan kondisi pengurangan pada langkah itu adalah
\[
  \kappa=\frac{|b|+\sqrt{b^{2}-4ac}}{\left|\sqrt{b^{2}-4ac}-|b|\right|}
  =\frac{2\times10^{8}}{1{,}49\times10^{-8}}\approx1{,}34\times10^{16},
\]
sehingga $\kappa u\approx1{,}49$: batas galat relatifnya melewati satu, artinya tidak
ada satu digit pun yang dijamin.

Gambar~\ref{fig:katastrofik} memisahkan dua hal yang sering tercampur pada langkah itu:
besar bilangannya, dan banyaknya digit yang masih berarti setelah pengurangan.

\begin{figure}[htbp]
\centering
\begin{tikzpicture}[every node/.style={font=\footnotesize}]
  \draw[ink,thick,->] (1.2,3.20) -- (10.6,3.20);
  \foreach \dec/\tk in {-9/{$10^{-9}$},-6/{$10^{-6}$},-3/{$10^{-3}$},0/{$10^{0}$},
                        3/{$10^{3}$},6/{$10^{6}$},9/{$10^{9}$}} {
    \draw[ink] ({1.4+(\dec+9)*0.5},3.26) -- ({1.4+(\dec+9)*0.5},3.14);
    \node[below] at ({1.4+(\dec+9)*0.5},3.12) {\tk};
  }
  \fill[accent] ({1.4+17*0.5},3.20) circle (2.4pt);
  \node[above=2pt] at ({1.4+17*0.5},3.24) {$b=10^{8}$};
  \fill[layerteal] ({1.4+1.173*0.5},3.20) circle (2.4pt);
  \node[above=2pt] at ({1.4+1.173*0.5},3.24) {$b-\sqrt{\Delta}$};
  \draw[<->,thick,dashed,layergrey] ({1.4+1.173*0.5},3.72) -- ({1.4+17*0.5},3.72);
  \node[above] at ({1.4+9.1*0.5},3.74) {enam belas orde besaran};
  \foreach \i in {0,...,16} {
    \fill[accent!28] ({1.4+\i*0.45},1.60) rectangle ++(0.45,0.40);
    \draw[ink!45] ({1.4+\i*0.45},1.60) rectangle ++(0.45,0.40);
  }
  \foreach \i in {0,...,15} {
    \fill[layergrey!16] ({1.4+\i*0.45},0.85) rectangle ++(0.45,0.40);
    \draw[ink!45] ({1.4+\i*0.45},0.85) rectangle ++(0.45,0.40);
  }
  \fill[accent!70] ({1.4+16*0.45},0.85) rectangle ++(0.45,0.40);
  \draw[ink!45] ({1.4+16*0.45},0.85) rectangle ++(0.45,0.40);
  \node[left] at (1.35,1.80) {$b$};
  \node[left] at (1.35,1.05) {$b-\sqrt{\Delta}$};
  \draw[{Latex}-{Latex},layergrey] (1.40,0.62) -- ({1.4+16*0.45},0.62);
  \node[below] at ({1.4+8*0.45},0.60) {enam belas digit sama, saling meniadakan};
  \node[right] at ({1.4+17*0.45},1.05) {satu digit};
\end{tikzpicture}
\caption{Atas: $b$ dan $\sqrt{b^{2}-4ac}$ pada skala log, berjarak enam belas orde
besaran dari selisihnya. Bawah: kotak adalah digit berarti, tujuh belas masuk dan enam
belas di antaranya sama sehingga saling meniadakan. Yang tersisa satu digit, yaitu
selisih $1{,}49\times10^{-8}$ itu sendiri; karena setengah ulp galat pada akar kuadrat
cukup untuk menggeser digit terakhir, akar kecil keluar dengan galat relatif $0{,}255$.}
\label{fig:katastrofik}
\end{figure}

Perbaikannya bukan menambah presisi, melainkan memilih akar yang dihitung lebih dulu.
Hasil kali kedua akar adalah $c/a$, jadi setelah akar besar $r_{1}$ diketahui, akar
kecil diperoleh dari pembagian, bukan dari pengurangan:
\[
  r_{1}=\frac{-b-\sqrt{b^{2}-4ac}}{2a},\qquad r_{2}=\frac{c}{a\,r_{1}}
  =\frac{2c}{-b-\sqrt{b^{2}-4ac}}.
\]
Bentuk kedua tidak pernah mengurangkan dua bilangan berdekatan. Untuk contoh di atas
ia memberi $-10^{-8}$ tepat, yaitu akar yang dibulatkan benar. Rumus yang sama dengan
tanda diatur menurut tanda $b$ dipakai di hampir setiap pustaka yang menyelesaikan
persamaan kuadrat, dan alasannya bukan kerapian melainkan kestabilan.

Wajah kedua pembatalan katastrofik muncul di tempat yang tampaknya tidak berbahaya:
selisih $1-\cos x$ untuk $x$ kecil. Untuk $x$ kecil,
\[
  \frac{1-\cos x}{x^{2}}\longrightarrow\frac{1}{2},
\]
karena $\cos x=1-x^{2}/2+x^{4}/24-\cdots$. Rumus itu benar, dan rumus itu juga yang
paling wajar ditulis. Tetapi $\cos x$ disimpan dengan galat relatif $u$, dan karena
$|\cos x|\le1$, galat mutlaknya sekitar $u$. Setelah dibagi $x^{2}$, galat itu menjadi
galat relatif sekitar $2u/x^{2}$. Untuk $x=10^{-8}$ faktor itu bernilai $2{,}2$:
seluruh digit lenyap.

Berikut hasil eksekusi nyata untuk sepuluh nilai $x$, dibandingkan dengan acuan yang
dihitung pada 50 digit dan dengan deret Taylor $0{,}5-x^{2}/24+x^{4}/720$; untuk
$x\le10^{-3}$ kedua pembanding itu sepakat sampai digit terakhir. Kolom ``lugu'' adalah
$(1-\cos x)/x^{2}$, kolom ``identitas'' adalah
$\tfrac12\left(\sin(x/2)/(x/2)\right)^{2}$ yang berasal dari
$1-\cos x=2\sin^{2}(x/2)$.

\begin{center}
\begin{tabular}{rccc}
\hline
$x$ & lugu & galat relatif & identitas \\
\hline
$10^{-1}$ & $0{,}499583472197418$ & $1{,}1\times10^{-14}$ & $0{,}499583472197423$ \\
$10^{-2}$ & $0{,}499995833347366$ & $2{,}9\times10^{-13}$ & $0{,}499995833347222$ \\
$10^{-3}$ & $0{,}499999958325503$ & $1{,}6\times10^{-11}$ & $0{,}499999958333335$ \\
$10^{-4}$ & $0{,}499999996961265$ & $5{,}2\times10^{-9}$ & $0{,}499999999583333$ \\
$10^{-5}$ & $0{,}500000041370185$ & $8{,}3\times10^{-8}$ & $0{,}499999999995833$ \\
$10^{-6}$ & $0{,}500044450291171$ & $8{,}9\times10^{-5}$ & $0{,}499999999999958$ \\
$10^{-7}$ & $0{,}499600361081320$ & $8{,}0\times10^{-4}$ & $0{,}500000000000000$ \\
$10^{-8}$ & $0{,}000000000000000$ & $1{,}0$ & $0{,}500000000000000$ \\
$10^{-9}$ & $0{,}000000000000000$ & $1{,}0$ & $0{,}500000000000000$ \\
$10^{-10}$ & $0{,}000000000000000$ & $1{,}0$ & $0{,}500000000000000$ \\
\hline
\end{tabular}
\end{center}

Polanya khas dan layak dihafal. Selama $x^{2}/2$ jauh di atas $u$, rumus lugu hampir
benar. Ketika $x^{2}/2$ menyusut ke skala $u$, digit berartinya habis dan galat
membesar mengikuti $2u/x^{2}$: setiap penurunan $x$ dengan faktor sepuluh menambah dua
digit kehilangan. Di $x=10^{-8}$, $1-\cos x$ lebih kecil daripada setengah ulp di
sekitar $1$, sehingga $\cos x$ dibulatkan tepat ke $1$ dan hasilnya nol. Titik runtuh
itu bukan misteri: $x^{2}/2\approx u$, yaitu $x\approx\sqrt{2u}\approx1{,}5\times10^{-8}$.

Kolom terakhir menunjukkan bahwa masalahnya berkondisi baik. Identitas sudut ganda
memberi jawaban benar sampai digit terakhir untuk setiap $x$ di tabel, karena ia tidak
pernah mengurangkan dua bilangan berdekatan. Galat bukan sifat masalah. Ia sifat rumus
yang dipilih.

\begin{resep}
Empat langkah memilih bentuk rumus yang aman dari pemusnahan:
\begin{enumerate}
  \item Tandai setiap pengurangan $a-b$ di dalam rumus, lalu hitung
        $\kappa=(|a|+|b|)/|a-b|$ pada data terburuk yang masuk akal.
  \item Terjemahkan $\kappa$ menjadi digit yang hilang: $\kappa u\approx10^{-d}$
        berarti $d$ digit terakhir tidak dapat dipercaya.
  \item Cari bentuk setara yang tidak memuat selisih itu: kalikan dengan sekawan,
        pakai identitas sudut ganda, faktorkan, atau panggil \texttt{log1p},
        \texttt{expm1}, dan \texttt{hypot}.
  \item Ulangi hitungan pada presisi lebih tinggi. Selisih kedua hasil harus sepadan
        dengan $\kappa u$; bila jauh lebih besar, bentuk barunya masih menyembunyikan
        pengurangan.
\end{enumerate}
Peringatan urutan: langkah 1 tidak boleh ditukar dengan langkah 3. Menulis ulang rumus
sebelum tahu pengurangan mana yang berbahaya hanya memindahkan masalah, seperti pada
bentuk satu-lintasan variansi yang menukar satu pengurangan berbahaya dengan
pengurangan berbahaya yang lain.
\end{resep}

\section{Penjumlahan: galat yang tumbuh bersama $n$}
\tingkat{lanjut}

Penjumlahan adalah operasi yang paling sering dianggap remeh, padahal ia tempat galat
paling mudah dihitung. Sebelum sampai ke batasnya, satu demonstrasi yang membongkar
anggapan bahwa urutan tidak penting:

\begin{center}
\texttt{(0.1 + 0.2) + 0.3 = 0.6000000000000001}, \quad
\texttt{0.1 + (0.2 + 0.3) = 0.6}.
\end{center}

Selisihnya $1{,}11\times10^{-16}$, tepat satu ulp pada skala $0{,}6$. Tidak ada
kesalahan tulis; hanya tiga bilangan yang masing-masing tidak punya wakil biner
berhingga, dan pembulatan yang berbeda urutan. Dalam matematika, penjumlahan bersifat
asosiatif. Di dalam mesin, ia asosiatif hanya bila kebetulan tidak ada pembulatan yang
berbeda.

\begin{teorema}[Galat penjumlahan berurutan]
Misalkan $s_{1}=x_{1}$ dan $s_{k}=\operatorname{fl}(s_{k-1}+x_{k})$ untuk
$k=2,\dots,n$ dengan $x_{i}\in\mathbb{F}$. Maka
\[
  \left|s_{n}-\sum_{i=1}^{n}x_{i}\right|
  \;\le\;(n-1)\,u\sum_{i=1}^{n}|x_{i}|+O(u^{2}).
\]
\end{teorema}

\begin{proof}
Setiap langkah memakai satu pembulatan, jadi $s_{k}=(s_{k-1}+x_{k})(1+\delta_{k})$
dengan $|\delta_{k}|\le u$. Menguraikan hasilnya secara rekursif memberi
\[
  s_{n}=\sum_{i=1}^{n}x_{i}\prod_{k=i}^{n}(1+\delta_{k}),
\]
dengan suku pertama tidak memuat $\delta_{1}$ karena $s_{1}=x_{1}$ belum dibulatkan.
Karena $\prod(1+\delta_{k})=1+\theta$ dengan
$|\theta|\le(1+u)^{n-1}-1\le(n-1)u+O(u^{2})$, ketaksamaan segitiga memberi batas yang
diklaim.
\end{proof}

Batas itu jujur tetapi longgar: ia menganggap setiap pembulatan mencapai galat
terburuknya dan semuanya searah, yang praktis tidak pernah terjadi. Tabel berikut
menunjukkan jurang antara batas dan kenyataan untuk deret harmonik
$H_{n}=\sum_{i=1}^{n}1/i$ dalam binary32, dengan nilai eksak dihitung oleh
\texttt{math.fsum} yang menjumlahkan secara eksak.

\begin{center}
\begin{tabular}{rccc}
\hline
$n$ & berurutan & eksak & batas $(n-1)u\sum|x_{i}|$ \\
\hline
$10^{2}$  & $5{,}1873779297$ & $5{,}1873775176$ & $3{,}1\times10^{-5}$ \\
$10^{3}$  & $7{,}4854784012$ & $7{,}4854708606$ & $4{,}5\times10^{-4}$ \\
$10^{4}$  & $9{,}7876129150$ & $9{,}7876060360$ & $5{,}8\times10^{-3}$ \\
$10^{5}$  & $12{,}0908508301$ & $12{,}0901461299$ & $7{,}2\times10^{-2}$ \\
$10^{6}$  & $14{,}3573579788$ & $14{,}3927267229$ & $8{,}6\times10^{-1}$ \\
\hline
\end{tabular}
\end{center}

Galat sesungguhnya pada $n=10^{6}$ adalah $3{,}54\times10^{-2}$, sekitar dua puluh lima
kali lebih kecil daripada batas, tetapi sudah menghapus dua setengah digit signifikan.
Galatnya juga tidak naik mulus: ia sempat turun pada $n=10^{4}$ karena pembulatan
searah dan berlawanan kebetulan saling meniadakan. Yang bisa diandalkan bukan besar
galat pada satu $n$, melainkan kecenderungannya membesar, dengan laju terburuk
sebanding $n$.

Penyebabnya bukan pembulatan yang besar, melainkan akumulator yang menjauh dari nol.
Begitu $s$ bernilai sekitar $14$, suku $1/i$ untuk $i$ besar berada di bawah setengah
ulp $s$ dan setiap penambahan kehilangan seluruh informasinya. Kasus ekstremnya mudah
dibuat: jumlahkan $1$ dengan sejuta suku bernilai $10^{-16}$; secara eksak hasilnya
$1{,}0000000001$, sedangkan penjumlahan berurutan tetap $1$ untuk berapa pun banyaknya
suku. Yang hilang bukan satu suku, melainkan seluruh kontribusi.

Gambar~\ref{fig:galatjumlah} merangkum pola yang sama untuk kelima nilai $n$ pada tabel
tadi, kali ini dengan penjumlahan Kahan sebagai pembanding.

\begin{figure}[htbp]
\centering
\begin{tikzpicture}
\begin{axis}[
  width=0.88\linewidth, height=6.0cm,
  xmode=log, ymode=log,
  xlabel={$n$}, ylabel={$|s_{n}-H_{n}|$},
  xmin=70, xmax=1.6e6, ymin=1e-8, ymax=1,
  grid=both, grid style={layergrey!15},
  tick label style={font=\footnotesize},
  label style={font=\footnotesize},
  legend style={font=\footnotesize,draw=layergrey!40,
                at={(0.03,0.97)},anchor=north west},
  legend cell align=left,
]
\addplot[ink,thick,mark=*,mark size=1.7pt] coordinates {
  (100,4.12e-7) (1000,7.54e-6) (10000,6.88e-6) (100000,7.05e-4) (1000000,3.54e-2)};
\addlegendentry{berurutan}
\addplot[accent,thick,mark=square*,mark size=1.6pt] coordinates {
  (100,6.48e-8) (1000,8.88e-8) (10000,2.03e-7) (100000,6.51e-8) (1000000,1.75e-7)};
\addlegendentry{Kahan}
\end{axis}
\end{tikzpicture}
\caption{Galat mutlak jumlah harmonik dalam binary32 pada lima ukuran $n$, skala
ganda-log. Penjumlahan berurutan naik dari $4{,}1\times10^{-7}$ pada $n=10^{2}$ ke
$3{,}54\times10^{-2}$ pada $n=10^{6}$; Kahan bertahan di orde $10^{-7}$ sepanjang
rentang, sehingga pada $n=10^{6}$ keduanya terpaut $2\times10^{5}$ kali. Data, presisi,
dan banyaknya operasi sama; yang berbeda hanya urutan penjumlahannya.}
\label{fig:galatjumlah}
\end{figure}

\begin{tangan}
Bahan: sebelas kartu kosong, spidol, dan satu aturan main. Setiap kali dua kartu
dijumlahkan, tulis hasilnya pada kartu baru dengan paling banyak empat angka penting,
lalu buang kartu lamanya.

Tulis $1{,}000$ pada satu kartu dan $0{,}0001$ pada sepuluh kartu lain. Jalur pertama:
letakkan kartu besar, lalu tambahkan sepuluh kartu kecil satu per satu. Setiap kali
$1{,}000+0{,}0001=1{,}0001$ dipotong menjadi $1{,}000$, sehingga kartu kecil itu hilang
tanpa jejak dan tumpukan tetap $1{,}000$.

Kocok ulang kartunya, jalur kedua: jumlahkan dulu sepuluh kartu kecil menjadi
$0{,}0010$, baru tambahkan kartu besar. Tumpukan akhir $1{,}001$.

Yang harus terlihat: sebelas kartu yang sama, aturan pemotongan yang sama, dua hasil
yang berbeda pada digit terakhir. Yang berubah hanya urutan, dan mesin mengulang urutan
itu sejuta kali.
\end{tangan}

Penjumlahan Kahan menyerang masalah itu secara langsung.\index{penjumlahan Kahan}
Idenya sederhana: setiap kali
sebuah suku hilang karena pembulatan, simpan sisa yang hilang di variabel kompensasi,
lalu masukkan kembali pada langkah berikutnya.

\begin{quote}
\texttt{s = 0.0; c = 0.0}\\
\texttt{for x in data:}\\
\hspace*{1.5em}\texttt{y = x - c}\\
\hspace*{1.5em}\texttt{t = s + y}\\
\hspace*{1.5em}\texttt{c = (t - s) - y}\\
\hspace*{1.5em}\texttt{s = t}
\end{quote}

Baris \texttt{c = (t - s) - y} adalah intinya, dan ia bekerja justru karena Sterbenz:
$t-s$ adalah selisih dua bilangan berdekatan sehingga eksak, dan $(t-s)-y$ karena itu
tepat sama dengan galat yang baru saja dibuang. Hasilnya tidak lagi menumpuk galat
sebanding $n$.

\begin{teorema}[Batas penjumlahan Kahan]
Dengan algoritme Kahan di atas, jumlah yang terhitung memenuhi ketaksamaan
\[
  s=\sum_{i=1}^{n}x_{i}(1+\delta_{i})+O(nu^{2})\sum_{i=1}^{n}|x_{i}|,
  \qquad |\delta_{i}|\le2u .
\]
\end{teorema}

Bukti teorema ini tidak kita ulang di sini. Ia adalah Teorema 8 pada kitab Goldberg,
dan buktinya ada di lampiran kitab itu, diambil dari Knuth. Yang perlu dibaca dari
pernyataannya: galat relatif tiap suku dibatasi $2u$, dua kali lipat galat satu
pembulatan, tetapi faktor $n$ hilang dari suku utama. Untuk $n$ kecil batas ini justru
lebih longgar daripada batas lugu, karena $2u>u$. Keunggulannya baru muncul ketika
$n$ besar, dan di situlah seluruh gunanya.

Angka nyatanya pada deret harmonik dengan $n=10^{6}$: penjumlahan berurutan dalam
binary32 memberi galat $3{,}54\times10^{-2}$, sedangkan Kahan memberi
$1{,}75\times10^{-7}$, yaitu $0{,}18$ ulp dari jawaban benar. Dalam binary64
selisihnya sama tajam: berurutan meleset $7{,}35\times10^{-13}$, atau $414$ ulp,
sedangkan Kahan menghasilkan angka yang sama persis dengan \texttt{math.fsum}.

Dua perbaikan lain layak disebut. Penjumlahan berpasangan, yang dipakai
\texttt{numpy.sum},\index{penjumlahan berpasangan} membagi data secara rekursif
sehingga galatnya tumbuh seperti $\log n$; ia hampir gratis dan cocok untuk data di
memori. \texttt{math.fsum} memakai akumulator eksak dan memberi hasil yang dibulatkan
benar untuk seluruh daftar, dengan biaya waktu lebih tinggi. Menjumlahkan dari suku
terkecil ke terbesar adalah perbaikan termurah, tetapi ia hanya mengurangi galat.

\section{Panduan praktis}
\tingkat{dasar}

Semua yang sudah kita bangun bisa diringkas menjadi beberapa kebiasaan yang murah
dilakukan dan mahal kalau diabaikan.

Pertama, pilih presisi berdasarkan bilangan kondisi, bukan berdasarkan kebiasaan.
Galat relatif hasil kira-kira $\kappa u$. Dengan binary64, sebuah perhitungan
berkondisi $10^{6}$ masih menyisakan sepuluh digit; dengan binary32, perhitungan yang
sama hanya menyisakan satu digit. Untuk $10^{8}$ operasi berurutan batas kasarnya
$nu$, dan dalam binary32 itu sudah melewati satu, artinya jawabannya tidak bermakna.

Kedua, jangan menguji kesamaan bilangan pecahan dengan \texttt{==}. Dua lintasan yang
secara matematis identik boleh menghasilkan bit yang berbeda, seperti
$0{,}1+0{,}2$ tadi. Uji yang benar memakai toleransi berskala, misalnya
$|x-y|\le10^{-12}\max(1,|x|,|y|)$, dan toleransi itu harus berasal dari analisis
galat, bukan dari percobaan sampai uji lulus.

Ketiga, curigai setiap rumus yang memuat selisih dua besaran hampir sama. Bentuk baku
variansi, $(1-\cos x)/x^{2}$, dan $\log(1+x)$ untuk $x$ kecil semuanya bisa ditulis
ulang. Pustaka standar sudah melakukannya: \texttt{log1p}, \texttt{expm1}, dan
\texttt{hypot} ada justru karena bentuk lugunya tidak stabil.

Keempat, untuk jumlah yang panjang atau bersuku sangat berbeda besarnya, gunakan
\texttt{math.fsum} atau penjumlahan berpasangan, dan simpan akumulator dalam presisi
lebih tinggi bila tersedia. Urutkan dari suku terkecil bila tidak ada pilihan lain.

Kelima, jangan menyimpulkan kebenaran dari kesepakatan angka. Dua program yang sepakat
sampai dua belas digit cukup meyakinkan; yang sepakat sampai tiga digit hanya memberi
tahu bahwa keduanya memakai rumus yang sama. Uji yang lebih baik adalah mengulang
hitungan pada presisi berbeda: bila hasilnya bergeser jauh melebihi ramalan
$\kappa u$, ada yang salah dalam rumusnya, bukan dalam mesinnya.


\pagebreak[4]
\begin{tantangan}
Hitung $10^{8}\left(\sqrt{10^{16}+1}-10^{8}\right)$ dalam binary64. Nilai sebenarnya
$0{,}5$, tetapi program yang menuliskan rumus itu apa adanya melaporkan $0{,}0$.

Jelaskan mengapa hasilnya nol, lalu hitung bentuk yang memberi $0{,}5$ tanpa menaikkan
presisi sama sekali.

Petunjuk 1: tanyakan berapa $\operatorname{ulp}(10^{16})$, lalu apakah $10^{16}+1$
masih punya tempat di grid bilangan mesin.

Petunjuk 2: kalikan pembilang dan penyebut dengan sekawan
$\sqrt{10^{16}+1}+10^{8}$; pengurangan berbahaya itu menjadi pembagian.

Kalau mentok, kembali ke bagian Pembatalan katastrofik: dua wajah.
\end{tantangan}

\section{Contoh hitung}
\tingkat{menengah}

Semua angka di bagian ini bisa diperiksa dengan kalkulator, \texttt{struct}, dan
pecahan biasa. Kita mulai dari membongkar satu pola bit, yaitu $0{,}1$ dalam binary32:

\[
  \underbrace{0}_{1}\;\underbrace{01111011}_{8}\;
  \underbrace{10011001100110011001101}_{23}.
\]

Eksponen tersimpan $123$, jadi setelah dikurangi bias $127$ diperoleh $e=-4$. Bilangan
bulat 23 bit pada mantissa adalah $5033165$, sehingga signifikannya
$1+5033165/2^{23}$. Maka
\[
  x=\left(1+\frac{5033165}{2^{23}}\right)2^{-4}=\frac{13421773}{134217728}.
\]
Pembilangnya berasal dari $2^{23}+5033165=13421773$. Sementara itu penyebutnya
$134217728$, dan sepersepuluhnya $13421772{,}8$, sehingga galatnya tepat
$0{,}2/134217728$, yaitu $1{,}4901161\times10^{-9}$,
dan galat relatifnya $1{,}49\times10^{-8}$, yaitu $u_{32}/4$. Angka itu masuk akal:
bilangan $1/10$ sama dengan $1{,}1001100110011\dots_{2}\times2^{-4}$, yang memerlukan
mantissa periodik dan tidak berhenti, sehingga pembulatan benar memotongnya pada bit
ke-24. Pola yang tersimpan
berakhir \texttt{\dots 1101}, bukan \texttt{\dots 1100}, karena sisa di belakangnya
lebih besar daripada setengah ulp dan pembulatan naik.

\kertas{Bongkar pola bit \texttt{0 01111100 01000000000000000000000} dengan tangan:
tentukan tanda, eksponen setelah dikurangi bias, nilai mantissa, lalu tulis bilangannya
sebagai pecahan dan sebagai desimal. Hitung $\operatorname{ulp}$ bilangan itu, lalu
putuskan hasil dari menambahkan $2^{-27}$ dan $3\cdot2^{-27}$ dengan aturan
\emph{tie to even}. Periksa jawabanmu dengan \texttt{struct.unpack} dan
\texttt{numpy.float32}. Bilangan ini dipilih agar mantissanya berhenti setelah dua
digit, sehingga seluruh aritmetikanya bisa dikerjakan tanpa kalkulator.}

Sekarang kasus $0{,}1+0{,}2$ dalam binary64, yang seluruhnya bisa dihitung dengan
pecahan. Ketiga bilangan yang terlibat berpenyebut pangkat dua:
\[
  \operatorname{fl}(0{,}1)=\frac{3602879701896397}{2^{55}},\qquad
  \operatorname{fl}(0{,}2)=\frac{3602879701896397}{2^{54}},
\]
\[
  \operatorname{fl}(0{,}3)=\frac{5404319552844595}{2^{54}}.
\]
Jumlah dua yang pertama sebelum dibulatkan adalah $10808639105689191/2^{55}$,
sedangkan $\operatorname{fl}(0{,}3)=10808639105689190/2^{55}$. Jadi jumlah eksak itu
terletak tepat setengah ulp di atas $\operatorname{fl}(0{,}3)$: titik seri, bukan di
dekatnya. Digit terakhir mantissanya $1$, yaitu ganjil, sehingga \emph{tie to even}
membulatkan ke atas: $0{,}300000000000000044408920985006\dots$

Selisihnya terhadap $\operatorname{fl}(0{,}3)$ karena itu tepat satu ulp,
$2^{-54}=5{,}551115123125783\times10^{-17}$, bukan galat acak. Pengurangan sesudahnya
eksak menurut Sterbenz, jadi angka yang tercetak program adalah angka itu persis.

Angka yang sama menjelaskan mengapa kesimpulannya bergantung pada presisi: dalam
binary32, $0{,}1+0{,}2$ dan $0{,}3$ menghasilkan pola bit yang identik karena kedua
galat pembulatannya kebetulan saling meniadakan. Jadi \texttt{0.1+0.2 != 0.3} bukan
hukum alam, melainkan akibat kebetulan pada satu presisi.

Sekarang akar kuadrat tadi, dengan aritmetika ulp. Diskriminan $b^{2}-4ac$ tepat
sama dengan $10^{16}-4$, karena $10^{16}$ masih di bawah $2^{54}$ dan bilangan itu
kelipatan $4$. Akar kuadratnya benar sampai $0{,}46$ ulp, jadi ia jatuh satu ulp di
bawah $b$ karena $\operatorname{ulp}(b)=2^{-26}$. Selisih sebenarnya $2\times10^{-8}=1{,}34$ ulp,
selisih yang terhitung $1$ ulp; kesalahan $0{,}34$ ulp pada skala $10^{8}$ itu, dibagi
dua, menjadi seluruh galat akar kecil:
\[
  r_{2}^{\text{naif}}=-\frac{2^{-26}}{2}=-2^{-27}
  =-7{,}450580596923828\times10^{-9},
  \qquad r_{2}=-10^{-8}.
\]
Bentuk stabil $2c/(-b-\sqrt{b^{2}-4ac})$ tidak mengalami pengurangan berbahaya, dan
hasilnya $-10^{-8}$ tepat, akar yang dibulatkan benar.

\section{Lab komputasi}
\tingkat{lanjut}

\begin{lab}
\begin{verbatim}
import math, struct
import numpy as np

def bit32(x):
    return ''.join(f'{b:08b}' for b in struct.pack('>f', float(x)))

print("0.1     :", bit32(0.1))
print("6.0     :", bit32(6.0))
print("-12.375 :", bit32(-12.375))
print("hex     :", (0.1).hex(), (6.0).hex(), (1e-8).hex())
print("eps     :", np.finfo(np.float32).eps, np.finfo(np.float64).eps)
print("ulp     :", math.ulp(1.0), math.ulp(1e10))
print("seri    :", 1 + 2.0**-53, 1 + 3*2.0**-53)

s = (0.1 + 0.2) - 0.3
print("0.1+0.2-0.3 :", repr(s), "=", s / math.ulp(0.3), "ulp")
print("asosiatif   :", repr((0.1 + 0.2) + 0.3), repr(0.1 + (0.2 + 0.3)))

n = 10**6
exact = math.fsum(1.0 / i for i in range(1, n + 1))

def jumlah(prec, kahan):
    s, c = prec(0), prec(0)
    for i in range(1, n + 1):
        y = prec(1) / prec(i)
        if kahan:
            y = y - c
            t = s + y
            c = (t - s) - y
            s = t
        else:
            s = s + y
    return s

print("H_n eksak   :", f"{exact:.12f}")
for prec, nama in ((np.float32, "float32"), (np.float64, "float64")):
    for kahan in (False, True):
        v = float(jumlah(prec, kahan))
        print(f"{nama} {'Kahan' if kahan else 'naif ':7} {v:.12f}"
              f"  galat {abs(v - exact):.2e}")

print(f"{'x':>8} {'lugu':>19} {'galat':>9} {'identitas':>19} {'galat':>9}")
for k in range(1, 11):
    x = 10.0**-k
    lugu = (1 - math.cos(x)) / x**2
    stabil = 2 * math.sin(x / 2)**2 / x**2
    acuan = 0.5 - x**2/24 + x**4/720 - x**6/40320
    print(f"{x:8.0e} {lugu:19.15f} {abs(lugu-acuan)/acuan:9.1e}"
          f" {stabil:19.15f} {abs(stabil-acuan)/acuan:9.1e}")

a, b, c = 1.0, 1e8, 1.0
d = b*b - 4*a*c
sq = math.sqrt(d)
print("diskriminan :", d, " akar:", sq)
print("b - akar    :", b - sq, " ulp(b):", math.ulp(b))
print("akar naif   :", (-b + sq) / (2*a))
print("akar stabil :", 2*c / (-b - sq))
\end{verbatim}
Keluaran sebenarnya:
\begin{verbatim}
0.1     : 00111101110011001100110011001101
6.0     : 01000000110000000000000000000000
-12.375 : 11000001010001100000000000000000
hex     : 0x1.999999999999ap-4 0x1.8000000000000p+2 0x1.5798ee2308c3ap-27
eps     : 1.1920929e-07 2.220446049250313e-16
ulp     : 2.220446049250313e-16 1.9073486328125e-06
seri    : 1.0 1.0000000000000004
0.1+0.2-0.3 : 5.551115123125783e-17 = 1.0 ulp
asosiatif   : 0.6000000000000001 0.6
H_n eksak   : 14.392726722866
float32 naif    14.357357978821  galat 3.54e-02
float32 Kahan   14.392726898193  galat 1.75e-07
float64 naif    14.392726722865  galat 7.35e-13
float64 Kahan   14.392726722866  galat 0.00e+00
       x                lugu     galat           identitas     galat
   1e-01   0.499583472197418   5.6e-15   0.499583472197423   5.4e-15
   1e-02   0.499995833347366   2.9e-13   0.499995833347222   1.1e-16
   1e-03   0.499999958325503   1.6e-11   0.499999958333335   0.0e+00
   1e-04   0.499999996961265   5.2e-09   0.499999999583333   1.1e-16
   1e-05   0.500000041370185   8.3e-08   0.499999999995833   1.1e-16
   1e-06   0.500044450291171   8.9e-05   0.499999999999958   0.0e+00
   1e-07   0.499600361081320   8.0e-04   0.500000000000000   1.1e-16
   1e-08   0.000000000000000   1.0e+00   0.500000000000000   0.0e+00
   1e-09   0.000000000000000   1.0e+00   0.500000000000000   0.0e+00
   1e-10   0.000000000000000   1.0e+00   0.500000000000000   0.0e+00
diskriminan : 9999999999999996.0  akar: 99999999.99999999
b - akar    : 1.4901161193847656e-08  ulp(b): 1.4901161193847656e-08
akar naif   : -7.450580596923828e-09
akar stabil : -1e-08
\end{verbatim}
\end{lab}

Tiga baris pertama memastikan tabel di bagian awal bab bukan karangan: pola bit
$0{,}1$, $6$, dan $-12{,}375$ keluar persis seperti hasil bongkar tangan, termasuk bit
tanda. Baris \texttt{hex} memperlihatkan tiga belas digit heksadesimal. Baris
\texttt{seri} adalah demonstrasi aturan genap yang paling telanjang: menambah
setengah ulp kepada $1$ tidak mengubah apa pun, menambah satu setengah ulp memberi
$1+2^{-51}$. Baris \texttt{0.1+0.2-0.3} mencetak
$5{,}551115123125783\times10^{-17}$ dan menyatakannya sendiri sebagai $1{,}0$ ulp;
baris \texttt{asosiatif} memberi dua hasil berbeda, berselisih
$1{,}11\times10^{-16}$.

Empat baris berikutnya adalah inti bab ini dalam bentuk paling ringkas: pada
$n=10^{6}$, binary32 berurutan kehilangan $3{,}54\times10^{-2}$ sedangkan Kahan hanya
$1{,}75\times10^{-7}$; binary64 berurutan $7{,}35\times10^{-13}$ sedangkan Kahan nol,
bit yang sama dengan \texttt{math.fsum}. Galat $7{,}35\times10^{-13}$ itu setara
$414$ ulp, jadi presisi ganda hanya memindahkan masalah ke skala lebih kecil.

Tabel $(1-\cos x)/x^{2}$ memperlihatkan pola yang sudah dihitung di kertas. Pada
$x=10^{-1}$ kedua cara berbeda dari acuan sekitar $5\times10^{-15}$ karena acuannya
sendiri terpotong pada suku $x^{6}$; mulai $x=10^{-2}$ acuan itu jauh lebih akurat
daripada mesin, dan galat rumus lugu naik dari $2{,}9\times10^{-13}$ ke $1{,}0$
ketika $x$ turun ke $10^{-8}$, sedangkan kolom identitas bertahan pada $10^{-16}$.
Nol pada $x=10^{-8}$ bukan pembulatan malas: $\cos(10^{-8})$ memang sama dengan $1$.

Blok terakhir memeriksa akar kuadrat: diskriminan $10^{16}-4$, akar $10^{8}-2^{-26}$,
selisih satu ulp $b$ yaitu $2^{-26}$, akar naif $-2^{-27}$, dan akar stabil $10^{-8}$
yang dibulatkan benar. Dua baris
terakhir berjarak $25{,}5$ persen, padahal rumusnya setara secara aljabar dan datanya
sama persis.

\begin{untukdosen}
Urutan mengajar yang berhasil dipakai: mulai dari satu kegagalan yang bisa dihitung
ulang, yaitu rumus kuadrat di awal bab, baru setelah itu masuk ke representasi. Bila
representasi datang lebih dulu, mahasiswa menghafal $1/8/23$ dan $1/11/52$ tanpa tahu
untuk apa. Setelah kegagalan itu, urutan yang bekerja adalah: banyaknya bilangan mesin
(agar jelas bahwa pembulatan tak terhindarkan), aturan pembulatan, lemma ulp, batas
$(1+\delta)$, lalu bilangan kondisi. Penjumlahan sebaiknya masuk setelah bilangan
kondisi, karena batas $(n-1)u\sum|x_{i}|$ baru terbaca sebagai ``galat yang tumbuh''
bila $\kappa$ sudah dikenal.

Dua miskonsepsi khas. Pertama, mahasiswa menganggap galat pembulatan bersifat acak dan
karena itu saling meniadakan pada penjumlahan panjang. Tabel harmonik membantahnya
dengan keras: galatnya bergerak searah, dan pada $n=10^{6}$ ia mencapai tiga puluh
lima ribu ulp. Kedua, mereka menganggap pengurangan selalu merusak digit, lalu
menghindari semua pengurangan. Teorema Sterbenz menunjukkan sebaliknya: pengurangan
dua bilangan mesin yang berdekatan justru eksak, dan yang merusak adalah galat yang
sudah ada pada masukannya. Aturan yang perlu diulang: yang berbahaya bukan tanda
minus, melainkan selisih kecil dari dua bilangan yang masing-masing sudah dibulatkan.

Aktivitas papan tulis yang murah dan efektif: minta seluruh kelas menghitung
$1+2^{-53}$ dan $1+3\cdot2^{-53}$ dengan aturan genap, lalu bandingkan dengan Python.
Setelah itu, minta dua sukarelawan mengerjakan
$(0{,}1+0{,}2)+0{,}3$ dan $0{,}1+(0{,}2+0{,}3)$ dengan menuliskan mantissa 52 bit
sebagai pecahan. Kegaduhan biasanya muncul saat mereka menemukan bahwa keduanya benar.

Tanda kesiapan lanjut: mahasiswa dapat menebak, tanpa menjalankan program, bahwa
$\operatorname{fl}(1+10^{-16})=1$, dapat menjelaskan mengapa bentuk
$2c/(-b-\sqrt{\Delta})$ lebih baik daripada $(-b+\sqrt{\Delta})/(2a)$, dan dapat
memperkirakan berapa digit yang hilang dari sebuah $\kappa$ yang diberikan. Bila
mereka menyebut kata ``galat ke belakang'' tanpa dipancing, mereka siap masuk Bab 25,
tempat gradien jaringan dihitung dengan aturan serupa.
\end{untukdosen}

\section{Ringkasan}
\tingkat{dasar}

\begin{ringkasan}
\begin{itemize}
  \item $\mathbb{F}$ berhingga: $2^{p-1}$ bilangan normal per pangkat dua, jadi
        hampir setiap bilangan real tidak punya wakil dan setiap penyimpanan berarti
        pembulatan.
  \item Pembulatan IEEE adalah \emph{round to nearest, ties to even}. Akibatnya
        $\operatorname{fl}(x)=x(1+\delta)$ dengan $|\delta|\le u=2^{-p}$, yaitu
        $1{,}11\times10^{-16}$ untuk binary64 dan $5{,}96\times10^{-8}$ untuk
        binary32.
  \item $\operatorname{eps}=2^{-p+1}$ adalah jarak $1$ ke bilangan sesudahnya;
        $\operatorname{ulp}(x)=2^{e-p+1}$ untuk $2^{e}\le|x|<2^{e+1}$. Galat satu
        pembulatan adalah setengah ulp, dan bila dinyatakan relatif ia berayun dengan
        faktor dua.
  \item Bilangan kondisi $\kappa_{f}=|xf'(x)|/|f(x)|$ mengukur pembesaran galat
        masukan. Untuk pengurangan, $\kappa=(|a|+|b|)/|a-b|$, yang membesar tanpa
        batas ketika $|a-b|$ mengecil.
  \item Pengurangan dua bilangan mesin yang berdekatan eksak (Sterbenz). Pembatalan
        katastrofik terjadi karena masukannya sudah membawa galat, bukan karena
        pengurangannya.
  \item Rumus kuadrat lugu menghitung akar kecil lewat pengurangan dua bilangan
        berukuran $|b|$; bentuk $2c/(-b-\sqrt{\Delta})$ menghindarinya dan memberi akar
        yang dibulatkan benar.
  \item Penjumlahan berurutan punya galat terburuk $(n-1)u\sum|x_{i}|$ dan tidak
        asosiatif. Kahan memindahkan galat ke variabel kompensasi sehingga faktor $n$
        hilang dari suku utama.
\end{itemize}
\end{ringkasan}

\section{Latihan}
\tingkat{menengah}

\begin{latihan}
\begin{enumerate}
  \item Tentukan $\operatorname{ulp}$ dari $1$; $1024$; $0{,}1$; dan $3\times10^{8}$
        dalam binary64, lalu bandingkan dengan $|x|\,u$. Berapa banyak bilangan
        binary64 yang terletak di $[1,2)$?

  \item Tunjukkan bahwa $1+2^{-53}=1$ dan $1+2^{-52}>1$ dalam binary64, lalu tentukan
        bilangan positif terkecil $x$ dengan $1+x>1$ dalam binary32. Nyatakan
        jawabanmu sebagai $2^{-k}$ dan bandingkan dengan tabel di bagian awal.

  \item Buktikan bahwa untuk setiap $x$ normal berlaku
        $|\operatorname{fl}(x)-x|\le\tfrac12\operatorname{ulp}(x)$, lalu jelaskan
        mengapa ketaksamaan itu tidak dapat diganti dengan kesamaan untuk semua $x$.

  \item Nilai klaim ini: \emph{karena setiap operasi IEEE-754 dibulatkan dengan benar,
        galat relatif hasil akhir program yang hanya memakai $+,-,\times,\div$ selalu
        paling besar beberapa $u$.} Putuskan benar, salah, atau belum cukup bukti, dan
        sebutkan satu contoh dari bab ini yang mendukung keputusanmu.

  \item Hitung dengan tangan $(0{,}1+0{,}2)+0{,}3$ dan $0{,}1+(0{,}2+0{,}3)$ sebagai
        pecahan berpenyebut pangkat dua, lalu tunjukkan bahwa selisihnya satu ulp.
        Periksa dengan \texttt{math.fsum} dan dengan \texttt{float.hex()}.
        (Petunjuk: nyatakan setiap bilangan sebagai $m/2^{55}$ atau $m/2^{54}$ dan
        hitung pembilangnya sebagai bilangan bulat.)

  \item Ulangi perhitungan $(1-\cos x)/x^{2}$ untuk $x=10^{-7}$ dengan identitas sudut
        ganda, lalu bandingkan dengan $0{,}5-x^{2}/24+x^{4}/720$. Berapa galat relatif
        sisanya, dan dari mana asalnya?

  \item Untuk $a=1$, $b=10^{4}$, $c=1$, hitung kedua akar dengan rumus lugu dan dengan
        bentuk stabil dalam binary64. Tentukan bilangan kondisi pengurangannya, lalu
        ramalkan berapa digit yang hilang sebelum menjalankan program. Periksa
        ramalanmu dengan Python.
        (Petunjuk: $\sqrt{10^{8}-4}\approx10^{4}-2\times10^{-4}$; bandingkan dengan
        $\operatorname{ulp}(10^{4})=2^{-39}$.)

  \item Buktikan bahwa penjumlahan berpasangan atas $n=2^{k}$ suku mempunyai galat
        terburuk sebanding $ku\sum|x_{i}|$, bukan $nu\sum|x_{i}|$, dengan menguraikan
        pohon penjumlahannya dan menghitung kedalaman pohon.

  \item Misalkan $x_{i}=10^{-16}$ untuk $i=1,\dots,10^{6}$ dan $x_{0}=1$. Hitung
        penjumlahan berurutan dan penjumlahan Kahan keduanya dengan tangan (boleh
        dibantu kalkulator), lalu jelaskan mengapa jawaban Kahan berbeda padahal
        setiap langkahnya juga dibulatkan.
        (Petunjuk: pada penjumlahan berurutan, setiap suku di bawah setengah ulp
        akumulator; pada Kahan, kompensasi $c$ menyimpan sisa itu dan menambahkannya
        kembali pada langkah berikutnya.)

  \item Sebuah program menghitung variansi dengan
        $\frac1n\sum x_{i}^{2}-\left(\frac1n\sum x_{i}\right)^{2}$ pada data yang
        seluruhnya sekitar $999$. Tentukan bilangan kondisi bentuk itu, ramalkan
        berapa digit yang hilang dalam binary32, dan bandingkan dengan bentuk
        dua-lintasan $\frac1n\sum(x_{i}-\bar{x})^{2}$. Jelaskan mengapa bentuk kedua
        lebih baik meskipun memerlukan dua kali lintasan data.
\end{enumerate}
\end{latihan}

\begin{untukpeneliti}
Yang benar-benar terbukti di bab ini adalah pernyataan-pernyataan baku dan tidak
bergantung pada mesin tertentu: batas $\operatorname{fl}(x)=x(1+\delta)$ dengan
$|\delta|\le u$, lemma ulp, teorema Sterbenz, batas galat penjumlahan berurutan, dan
ketaksamaan bilangan kondisi lewat teorema nilai rata-rata. Rujukan primernya Goldberg,
\emph{What Every Computer Scientist Should Know About Floating-Point Arithmetic}
(ACM Computing Surveys 23(1), 1991), bagian ``Relative Error and Ulps'',
``Cancellation'', ``Exactly Rounded Operations'', dan ``Errors in Summation''. Batas
Kahan adalah Teorema 8 kitab itu; buktinya di lampiran, diambil dari Knuth, dan tidak
diulang di sini. Untuk perlakuan sistematis mengenai kestabilan, rujukan lanjutannya
Higham, \emph{Accuracy and Stability of Numerical Algorithms}, bab tentang analisis
galat dan penjumlahan.

Yang lebih sering dipopulerkan daripada dibuktikan ada tiga. Pertama, klaim bahwa
``galat pembulatan bersifat acak dan saling meniadakan''. Itu benar hanya sebagai
model statistik untuk pembulatan yang tak berkorelasi, dan justru gagal pada kasus
yang paling penting: akumulasi pada penjumlahan panjang bersifat sistematis, seperti
terlihat pada tabel harmonik. Kedua, klaim bahwa galat selalu bisa dikurangi dengan
presisi lebih tinggi. Untuk pembatalan katastrofik, menaikkan presisi hanya
memindahkan titik runtuh, bukan menghapusnya: dalam binary64 titik itu ada di
$x\approx1{,}5\times10^{-8}$, dan dalam binary128 ia ada di sekitar $10^{-17}$. Ketiga,
klaim bahwa ulp adalah satuan galat yang bebas ambiguitas. Goldberg menunjukkan bahwa
ulp dapat didefinisikan sebagai jarak antar bilangan atau setengah jarak itu, dan
selisih kedua definisi itu muncul sebagai faktor dua dalam setiap batas. Angka
``$0{,}46$ ulp'' pada pembahasan akar kuadrat di bab ini memakai definisi pertama.

Batas keberlakuan hasil di sini perlu dinyatakan. Seluruh analisis menganggap model
pembulatan benar dan tidak ada \emph{fused multiply-add} yang mengubah urutan operasi,
tidak ada pengaturan ulang oleh pengompilasi, dan tidak ada mode pembulatan selain
\emph{to nearest}. Di perangkat nyata, vektorisasi dan instruksi FMA dapat mengubah
hasil bit demi bit; itu bukan pelanggaran standar, tetapi membuat perbandingan antar
mesin menjadi tidak deterministik. Analisis galat juga tidak menyentuh galat model,
galat data, atau galat pemotongan: bab ini hanya membahas galat pembulatan.

Pertanyaan terbuka yang jujur. Pertama, analisis galat klasik memberi batas terburuk,
sedangkan yang dibutuhkan praktisi adalah perkiraan yang ketat untuk kasus khas;
probabilistik pembulatan masih menjadi alat yang aktif diteliti dan belum menjadi
teori yang selesai. Kedua, untuk algoritme iteratif dan penyelesaian sistem linear,
galat pembulatan berinteraksi dengan bilangan kondisi matriks sedemikian rupa sehingga
batas naif terlalu pesimistis; analisis ke belakang, yang menjadi gagasan utama
Higham dan Wilkinson, memberi kerangka yang lebih tajam tetapi tidak dibahas di sini.
Ketiga, keputusan praktis seperti ``apakah binary32 cukup'' bergantung pada distribusi
data, dan tidak ada jawaban yang berlaku umum.
\end{untukpeneliti}

\peta{Bab ini mengikuti Goldberg, \emph{What Every Computer Scientist Should Know About
Floating-Point Arithmetic} (ACM Computing Surveys, 1991), bagian ``Introduction'',
``Relative Error and Ulps'', ``Cancellation'', ``Exactly Rounded Operations'',
``Formats and Operations'', ``Special Quantities'', dan ``Errors in Summation'', dengan
Teorema 8 untuk batas penjumlahan Kahan. Untuk sisi analisis real dari gagasan limit
dan teorema nilai rata-rata yang dipakai pada bukti bilangan kondisi, lihat Abbott,
\emph{Understanding Analysis}, bagian 2.2 (Limit of a Sequence), 5.3 (The Mean Value
Theorem), dan 6.6 (Taylor Series). Contoh $0{,}1+0{,}2$, pembongkaran pola bit,
analisis $(1-\cos x)/x^{2}$, dan seluruh angka pada tabel serta lab adalah perhitungan
sendiri yang dijalankan dengan Python 3.14 dan NumPy 2.5 dan diverifikasi ulang dengan
pecahan biasa serta aritmetika 50 digit. Untuk perlakuan lanjutan mengenai kestabilan
dan galat ke belakang, rujukan yang biasa dipakai adalah Higham, \emph{Accuracy and
Stability of Numerical Algorithms}.}
